\documentclass[openany,fontset=fandol]{ctexbook}
\usepackage{amsmath,amssymb}
\usepackage[a4paper,margin=25mm]{geometry}
\usepackage[unicode,hidelinks]{hyperref}
\setlength{\emergencystretch}{2em}
\raggedbottom
\newcommand{\PD}{\check P}
\newcommand{\Ys}{Y^{\mathrm{sing}}}
\newcommand{\Ds}{Y^{\mathrm{sup\,sing}}}
\begin{document}
\begin{center}
{\Large\bfseries 代数几何基础第五卷预备笔记}\par\medskip
{\LARGE\bfseries 双切超平面与特征零判据}\par\medskip
第107批：8.10--8.12\quad 2026年11月24日
\end{center}

\section*{底本与记号}
法稿物理43页的8.10起，续法44、45，止于8.13之前；对读英编稿物理30、31。承接8.8、8.9：$X/k$ 光滑、几何整且在 $P=\mathbb P^r_k$ 中局部闭浸入，正余维；$Z=\Ys$，$D=\Ds$，$T$ 是 $Z$ 在 $\PD$ 中像的约化闭包，$g:Z\to T$ 为诱导态射，$L=k(Z)$、$K=k(T)$。$U\subset T$ 为相对 $k$ 光滑处。像的余维均指其闭包的余维。

\section*{注8.10：无分歧但不 étale 的点}
在推论8.9的假设下，已有 $g^{-1}(U)\subset Z\setminus D$。即使 $X$ 在 $P$ 中闭，此包含也不必为等号。$g^{-1}(U)$ 恰是 $g$ 为 étale 的点集，$Z\setminus D$ 则是 $g$ 无分歧的点集。因此，可以有 $g$ 在其处无分歧而不 étale 的 $y\in Z$；这时 $g(y)$ 处的 $T$ 不是几何单支的，因而也不是几何正规的。\footnote{英30在两个点集之后误把差别归给整个 $Z\setminus D$，即把“无分歧”读成“无分歧但不 étale”。法43的原句是先定义两个点集，再说明差集可以非空。}

几何上，这对应于以下情形：某个超平面在 $x$ 处与 $X$ 相切而不密切，同时也在另一个点 $x'$ 处与 $X$ 相切。平面中次数至少四的非奇异曲线提供双切线的例子；双切线对应于对偶曲线上有两个局部分支的情形。\footnote{法44第二个切点处，打字末尾仍作 $x$，却紧接在“另一个点 $x'$”之后，故正文改为 $x'$。原稿另称“任意特征”，其双切线现象不可无条件当作8.9假设下的反例：特征二、$\dim X=1$ 时，极化型为一维交错型，只能为零，不能给出8.7(b)的普通奇异判据。这里正文叙述双切几何，例子的 étale 对比由下页特征零实例核实。双切线的存在也不自动断言对偶点恰为普通节点。}

\paragraph{译校补足：两个点集与几何单支。}
8.9给出 $g^{-1}(U)\subset\operatorname{Et}(g)$。反之，若 $g$ 在 $y$ 处 étale，则可在 $y$ 的一个开邻域上取其 étale 限制，像为包含 $g(y)$ 的开集；由此满射 étale 覆盖，源对 $k$ 的光滑下降到该像，故 $g(y)\in U$。又因 $T\hookrightarrow\PD$ 为闭浸入，$\Omega_{Z/T}\simeq\Omega_{Z/\PD}$，所以8.7(b)判出的正是 $g$ 的无分歧处。若 $T$ 在 $g(y)$ 几何单支，整源的泛点特化到 $y$ 且映到 $T$ 的泛点；由 \href{https://stacks.math.columbia.edu/tag/0GS8}{Stacks 0GS8}，此处无分歧会强制 étale，给出原断言。这里不需要假设 $g$ 在全部 $Z$ 上有限。

\newpage
\paragraph{译校例：特征零的严格包含。}
在代数闭的特征零域上，取平面四次曲线
\[
 F=(x^2-z^2)^2+yz^3+y^4=0,
 \qquad H:\ y=0.
\]
直线 $H$ 恰在 $[1:0:1]$、$[-1:0:1]$ 各有二重交；这两点 $F_y=1$，且切线方向的二次项非零，因而均非超奇异。曲线确实光滑：$z=0$ 时，$F_x=4x^3$、$F_y=4y^3$ 不会在射影点同时为零；$z=1$ 时 $F_x=4x(x^2-1)$，若三偏导都为零则 $x=0$ 或 $x=\pm1$。在 $x=0$ 时 $F_z=4+3y$ 强制 $y=-4/3$，于是 $F_y=1+4y^3\ne0$；在 $x=\pm1$ 时 $F_z=3y$ 强制 $y=0$，仍有 $F_y\ne0$。光滑平面曲线在代数闭域上整：否则其不同正次数分支依 Bézout 相交，交点不光滑。于是8.8的条件成立。$H$ 的纤维有两点；若 $H\in U$，由8.9的适当情形该纤维只有一点，矛盾。因此这两点都给出无分歧而不 étale 的实例。

\section*{推论8.11：特征零}
设 $k$ 的特征为零。
\begin{enumerate}
\renewcommand{\labelenumi}{(\alph{enumi})}
\item $D$ 在 $\PD$ 中的像余维至少为二。
\item 8.8的条件 (iii) 与其余条件等价。换言之，其余条件的否定 $Z=D$（就底集而言），也意味着 $Z$ 在 $\PD$ 中的像余维至少为二。
\end{enumerate}
结合8.8，由 (b) 推出 (a)。而由维数理论，(iii) 等价于 $L/K$ 为有限扩张（可把这写成8.8的 (iii bis)）；在特征零下它自动可分，于是得到8.8的 (i bis)。\footnote{法44在(b)末回到奇异零集的像，即 $Z$；如果写作 $D$ 就退回(a)，不能表达(b)对全部奇异参数的结论。英30在此并列了两种转写。}

\paragraph{译校补足：对 (b) 推出 (a) 分情形。}
若 (iii) 成立，则8.8已经证明 $D$ 的像闭包余维至少为二。若 (iii) 不成立，则 $\dim T<r-1$，而 $D$ 的像包含于 $T$，仍有余维至少二。这一步只用两种情形分别判断，不把“余维至少二”误读作“恰余维二”。维数判定来自
\[
 \operatorname{trdeg}_K L=\dim Z-\dim T=r-1-\dim T.
\]
函数域扩张为有限型，超越次数为零等价于有限扩张；在正特征下不能由有限性推出可分性。

\newpage
\section*{注8.12：一般线性笔}
\paragraph{(a)}
几何上，8.11(a)表示：取足够一般的超平面截面线性笔，每一个成员不是光滑，就是只有几何普通奇异点；事实上至多有一个这样的几何奇异点，因此可将(a)稍微加强。8.11(b)还表示，在特征零情形，如果 $Z=D$，那么足够一般的笔的所有成员都光滑。$Z=D$ 可用一个可逆模的行列式截面恒等于零来表达。\footnote{法44及英30可辨为 $\omega_{X/k}^{\otimes2}\otimes\mathcal O_Z(1,1)$ 一类写法，未见一般秩幂次。按二次型定义，令 $\mathcal L=\mathcal O_Z(1,1)$、$d=\dim X$，则相对极化型给出 $T_{X/k}|_Z\to\Omega^1_{X/k}|_Z\otimes\mathcal L$，其行列式属于 $\omega_{X/k}^{\otimes2}|_Z\otimes\mathcal L^{\otimes d}$。此处限制符号表示沿 $Z\to X$ 拉回。正文采用不歧义的行列式描述，不把可见的一次扭转当作一般公式。}

第二种情形，不论是否在特征零，都应视为很例外。原稿以试探口吻说，照经典术语，这可称 $X$ 对所考虑的射影浸入是“直纹的”；它建议用8.5及其推论说明这种术语，并说若 $\dim X=1$，这意味着 $X$ 在一条直线上。\footnote{法45这段保留“若无误”的措辞，旁有括线和未能完全辨认的手写 $T=L$ 等边注，英31也留下未辨字。译文不以边注补造普遍的分类定理。原稿“足够一般的笔全光滑”适用于全部奇异参数像余维至少二的情形，不能由任意正特征下的 $Z=D$ 单独推出。}

\paragraph{译校补足：为何同一截面至多一个普通奇异点。}
若8.8的等价条件成立，另须使笔避开 $T\setminus U$，而不只是避开 $D$ 的像。因为 $U$ 是 $T$ 的稠密光滑开集，$T\setminus U$ 在 $\PD$ 中余维至少二；8.9在 $U$ 上给出开浸入，每个几何纤维至多一点，且该点对应普通奇异零。若8.8的条件不成立而 $k$ 特征零，则可直接避开整个 $T$。余维至少二的闭集可由一般的几何直线避开，见第105批的直线入射维数论证。不把一般几何线性笔的存在解释为有限底域上必有有理笔。“光滑或只有普通奇异点”按奇异零判据处理，不把恒等为零的截面误计成期望维数的光滑截面。

\paragraph{(b)}
若特征为 $p>0$，按理应给出这样的光滑平面曲线例子：8.8的等价条件不成立，即 $Z=D$，但仍有 $\dim T=r-1$；亦即 $L/K$ 为有限不可分扩张。原稿说此处没有构造例子，却相信例子存在。\footnote{法45打字回指8.6，英31照录。按同句 $Z=D$ 和 $\dim T=r-1$，应分别对应8.8的 (ii) 的否定和 (iii)，正文因而改为8.8。}

\paragraph{译校例：正特征下的维数等式不推出泛 étale。}
取代数闭的特征 $p>0$ 底域上的 Fermat 曲线 $x^{p+1}+y^{p+1}+z^{p+1}=0$。三偏导为 $x^p,y^p,z^p$，无共同射影零点，所以曲线光滑。余法秩为一，故 $Z\simeq X$；切线参数为 $[x^p:y^p:z^p]$。这个映射经过相对 Frobenius，且其像仍是一条曲线（本例方程系数在素域）。在完美底域的曲线函数域上，$[L:L^p]=p$，故它诱导纯不可分的次数 $p$ 扩张，切映射为零，故处处超奇异，而 $\dim T=1=r-1$。

\smallskip
\noindent\textbf{原稿编排提示与接续。}
法45在下一命题前建议：在(a)中脚注指向后段，说明在例外“直纹”情形改变射影浸入可回到8.8的一般情形；还建议把8.6至此的微分内容单列，并把此前开头和下文并为讨论例外参数维数的一节。这是编排提议，非本批已证结论。下次由法45命题8.13标题接英31，须先辨大段划删除与手写替代，继读法46。本批未改1158页合订本。
\end{document}
